% !TeX root = ../tesi.tex

\capitolononnumerato{Glossario}

Raccolgo in ordine alfabetico i termini tecnici introdotti nel testo. Alla prima occorrenza nel
corpo della relazione, il termine e il pedice G rimandano alla voce corrispondente.

\begin{description}
\voceglossario{adr}{ADR}{
	Documento che registra una decisione di architettura: il contesto in cui nasce, le
	alternative considerate, la scelta e le sue conseguenze. Conserva le motivazioni della scelta anche
	quando cambia chi lavora al progetto.
}

\voceglossario{api}{API}{
	Interfaccia con cui un programma rende disponibili dati o operazioni ad altri programmi,
	secondo un contratto definito. In Shopify, la Admin API permette di amministrare il negozio; la
	Storefront API espone le funzioni rivolte all’esperienza d’acquisto.
}

\voceglossario{app-proxy}{\emph{App proxy}}{
	Funzione di Shopify che espone un percorso dell’applicazione sotto il dominio del negozio e
	inoltra le richieste al \emph{server} dell’applicazione aggiungendo una firma che ne attesta la
	provenienza. Permette a un \emph{widget} di chiamare il proprio \emph{server} senza uscire dal
	dominio del negozio.
}

\voceglossario{b-impact-assessment}{\emph{B Impact Assessment}}{
	Questionario gratuito con cui un’azienda misura il proprio impatto sociale e ambientale in cinque
	aree: \emph{governance}, lavoratori, clienti, comunità e ambiente. Il punteggio che ne risulta è
	anche la base per la certificazione B Corp.
}

\voceglossario{b2b}{B2B}{
	Modello commerciale in cui un’azienda vende prodotti o servizi ad altre aziende. Nel
	commercio elettronico può richiedere cataloghi, listini e condizioni di acquisto diversi per
	ciascun cliente.
}

\voceglossario{b2c}{B2C}{
	Modello commerciale in cui un’azienda vende prodotti o servizi al consumatore finale. È il
	contesto dell’interfaccia conversazionale di ChatCart, rivolta a chi acquista nel negozio.
}

\voceglossario{brain-dump}{\emph{Brain dump}}{
	Tecnica di organizzazione personale che consiste
	nello scrivere o registrare, senza ordine né filtro, tutto ciò che si ha in mente (attività,
	idee, dubbi), per poi ricavarne le azioni da svolgere. È il primo passo del metodo Getting
	Things Done di David Allen, che lo chiama raccolta.
}

\voceglossario{caso-uso}{Caso d’uso}{
	Descrizione di un’interazione completa tra un sistema e un
	attore, cioè un ruolo esterno che lo usa per raggiungere uno scopo: comprende lo scenario
	principale, in cui tutto va come previsto, e gli scenari alternativi, con le eccezioni e gli
	errori. I casi d’uso si rappresentano spesso con un diagramma UML, che mostra gli attori, il
	confine del sistema e i casi d’uso collegati a ciascun attore; da essi si ricavano i requisiti.
}

\voceglossario{chatbot}{\emph{Chatbot}}{
	Programma che conversa con l’utente in linguaggio naturale, per iscritto o a voce. Nei negozi
	\emph{online} risponde alle domande dei visitatori; quelli più recenti si basano su un modello
	linguistico e possono consultare i dati del negozio.
}

\voceglossario{commercio-agentico}{Commercio agentico}{
	Insieme delle esperienze d’acquisto in cui un agente di intelligenza artificiale accompagna il
	cliente, o agisce per suo conto, nella ricerca, nella scelta e nell’acquisto di un prodotto.
	L’agente può appartenere al negozio oppure a una terza parte che raggiunge il negozio tramite
	protocolli come MCP e UCP.
}

\voceglossario{conventional-commits}{Conventional Commits}{
	Convenzione per i messaggi di \emph{commit}
	di Git: ogni messaggio inizia con un tipo (\texttt{feat}, \texttt{fix}, \texttt{refactor},
	\texttt{docs}, \texttt{test}, \texttt{chore} e altri),
	eventualmente seguito da un ambito tra parentesi, e da una breve descrizione. Rende la storia
	leggibile da persone e programmi, per esempio per generare le note di rilascio o calcolare il
	numero di versione secondo SemVer.
}

\voceglossario{crm}{CRM}{
	Insieme di pratiche e sistemi per gestire le relazioni di un’azienda con i propri clienti.
	Un sistema CRM raccoglie, per esempio, contatti, interazioni commerciali e richieste di assistenza.
}

\voceglossario{dashboard}{\emph{Dashboard}}{
	Pagina che riassume in un’unica vista dati e indicatori utili a chi la consulta. La
	\emph{dashboard} di ChatCart è incorporata nell’amministrazione del negozio e mostra al
	commerciante il \emph{funnel}, le conversazioni e la loro analisi.
}

\voceglossario{design-pattern}{\emph{Design pattern}}{
	Soluzione riconosciuta e documentata a un problema di progettazione che si ripresenta in contesti
	diversi, descritta in modo indipendente dal linguaggio. Ne è un esempio l’\emph{adapter}, che
	rende compatibili due interfacce senza modificarle.
}

\voceglossario{design-system}{\emph{Design system}}{
	Insieme coerente di principi, regole visive e componenti riusabili per progettare
	un’interfaccia. Comprende elementi come colori, tipografia, spaziature e comportamenti dei
	componenti, così da mantenere riconoscibile l’identità del prodotto o del marchio.
}

\voceglossario{e-commerce}{\emph{E-commerce}}{
	Commercio di prodotti o servizi attraverso canali digitali. Un negozio elettronico comprende
	il catalogo, la selezione dei prodotti, il carrello e il pagamento, oltre ai collegamenti con i
	sistemi che gestiscono ordini e logistica.
}

\voceglossario{erp}{ERP}{
	Sistema gestionale che integra dati e processi di diverse aree aziendali, come
	amministrazione, acquisti, magazzino e vendite. Il collegamento con il negozio permette, per
	esempio, di sincronizzare ordini e disponibilità dei prodotti.
}

\voceglossario{founder-senior-consultant}{\emph{Founder \& Senior Consultant}}{
	Ruolo che unisce la fondazione dell’azienda all’attività di consulenza esperta. In Halue lo
	ricopre Valentino Baraldo.
}

\voceglossario{graphql}{GraphQL}{
	Linguaggio di interrogazione per API e insieme di regole per eseguire le richieste su uno
	schema di dati tipizzato. Il programma che invia la richiesta specifica i campi di cui ha bisogno;
	Shopify lo adotta per la Admin API e la Storefront API.
}

\voceglossario{headless}{\emph{Headless}}{
	Architettura che separa la gestione dei dati e delle funzioni applicative dall’interfaccia
	che li presenta. In un negozio di questo tipo l’interfaccia è un’applicazione autonoma che comunica
	con la piattaforma commerciale tramite API; in un CMS la stessa separazione riguarda i contenuti e
	la loro presentazione.
}

\voceglossario{hydrogen}{Hydrogen}{
	\emph{Framework} di Shopify basato su React per costruire interfacce di negozi
	\emph{headless}. Offre strumenti e componenti per accedere ai dati del negozio tramite la
	Storefront API; Oxygen è il servizio di Shopify che può ospitare queste applicazioni.
}

\voceglossario{interfaccia-generativa}{Interfaccia generativa}{
	Interfaccia in cui un modello linguistico decide, a ogni risposta, quali elementi mostrare
	all’utente oltre al testo. In ChatCart il modello sceglie tra componenti predefiniti e ne fornisce
	i dati; l’aspetto dei componenti resta fissato dal codice.
}

\voceglossario{ipaas}{iPaaS}{
	Piattaforma erogata come servizio per collegare applicazioni e sistemi diversi. Offre
	connettori e strumenti per trasformare e trasferire i dati, riducendo la necessità di costruire
	un’integrazione distinta per ogni coppia di sistemi.
}

\voceglossario{junior-developer}{\emph{Junior Developer}}{
	Sviluppatore all’inizio del proprio percorso professionale, che realizza e verifica
	funzionalità con il supporto di colleghi più esperti.
}

\voceglossario{landing-page}{\emph{Landing page}}{
	Pagina di atterraggio, costruita per accogliere chi arriva da una campagna, da un annuncio o da un
	collegamento esterno e guidarlo verso un’azione, come l’acquisto di un prodotto. Una \emph{landing
	page} dinamica cambia contenuti in base a dati come la collezione, il mercato o il visitatore.
}

\voceglossario{liquid}{Liquid}{
	Linguaggio di \emph{template} usato nei temi Shopify. Permette di combinare la struttura
	delle pagine con i dati del negozio, attraverso variabili, condizioni e cicli; la piattaforma
	elabora i modelli e genera le pagine da inviare al navigatore.
}

\voceglossario{llm}{LLM}{
	Modello linguistico di grandi dimensioni, addestrato su ampie raccolte di dati per elaborare
	e generare testo. In un agente può selezionare strumenti esterni di cui l’applicazione esegue le
	operazioni; il collegamento a questi strumenti gli permette di usare dati aggiornati che non fanno
	parte del suo addestramento.
}

\voceglossario{mcp}{MCP}{
	Protocollo aperto che standardizza la comunicazione tra applicazioni basate su modelli
	linguistici e servizi esterni. Un \emph{server} MCP può esporre strumenti, risorse e istruzioni; è
	l’applicazione che ospita il modello a gestire il collegamento e l’esecuzione delle richieste.
}

\voceglossario{oauth}{OAuth}{
	Standard aperto per l’autorizzazione delegata: permette a un servizio di concedere a
	un’applicazione un accesso limitato e revocabile ai dati di un utente, con \emph{token} di
	accesso al posto della password. Definisce più flussi (\emph{grant}), tra cui l’\emph{authorization
	code}, le \emph{client credentials} e il \emph{token exchange}, adatti a casi d’uso diversi.
}

\voceglossario{project-management-assistant}{\emph{Project Management Assistant}}{
	Figura che supporta la gestione dei progetti, contribuendo all’organizzazione delle
	attività, all’aggiornamento della documentazione e al monitoraggio delle scadenze.
}

\voceglossario{project-service-manager}{\emph{Project/service manager}}{
	Responsabilità operativa di gestione di un progetto o di un servizio continuativo. Nel
	metodo di Halue cura la pianificazione delle attività e la capacità disponibile del gruppo di
	lavoro.
}

\voceglossario{pull-request}{\emph{Pull request}}{
	Proposta di integrare un insieme di modifiche da un \emph{branch} a un altro in un sistema
	di versionamento. Raccoglie le differenze nel codice, la discussione di revisione e gli esiti dei
	controlli automatici prima dell’integrazione.
}

\voceglossario{redesign}{\emph{Redesign}}{
	Riprogettazione grafica e funzionale di un sito o di un negozio esistente, che ne rinnova
	l’interfaccia e i percorsi dell’utente, spesso senza cambiare la piattaforma su cui si basa.
}

\voceglossario{ricerca-semantica}{Ricerca semantica}{
	Ricerca che confronta il significato della richiesta con quello dei documenti, invece delle sole
	parole in comune. Di solito si basa su rappresentazioni numeriche dei testi, dette
	\emph{embedding}, in cui testi di significato simile risultano vicini.
}

\voceglossario{sal}{SAL}{
	Stato di avanzamento dei lavori. Nel contesto descritto indica un rapporto periodico che
	riassume il progresso del progetto, le scadenze, i rischi e i punti aperti, per condividere la
	situazione con le persone coinvolte.
}

\voceglossario{scrum}{Scrum}{
	\emph{Framework} per affrontare problemi complessi attraverso cicli brevi di lavoro,
	chiamati \emph{sprint}, e un confronto regolare sui risultati. Prevede responsabilità, eventi e
	artefatti definiti; nella \secref{sec:organizzazione-halue} distinguo questo riferimento
	dall’adattamento che Halue ne fa per il lavoro su commessa.
}

\voceglossario{semver}{SemVer}{
	Convenzione di versionamento semantico che esprime una versione nella forma
	\texttt{MAJOR.MINOR.PATCH}. Rispetto all’API pubblica dichiarata, il primo numero cambia per
	modifiche incompatibili, il secondo per nuove funzionalità compatibili e il terzo per correzioni
	compatibili.
}

\voceglossario{senior-project-manager}{\emph{Senior Project Manager}}{
	Figura con esperienza nella gestione dei progetti, responsabile del coordinamento di
	pianificazione, risorse e avanzamento delle attività.
}

\voceglossario{senior-software-engineer}{\emph{Senior Software Engineer}}{
	Professionista dello sviluppo del \emph{software} con esperienza consolidata, che affronta
	problemi progettuali complessi e contribuisce alle scelte tecniche e alla revisione del codice.
}

\voceglossario{shopify-markets}{Shopify Markets}{
	Funzione di Shopify per vendere in più paesi dallo stesso negozio. Per ogni mercato si
	definiscono lingua, valuta, prezzi e indirizzi delle pagine.
}

\voceglossario{shopify-plus}{Shopify Plus}{
	Proposta di Shopify per le aziende di grandi dimensioni, con funzioni aggiuntive rispetto ai piani
	standard, tra cui quelle per la vendita B2B.
}

\voceglossario{skill}{\emph{Skill}}{
	Insieme riusabile di istruzioni e, quando necessari, risorse o strumenti, che guida un
	assistente basato su IA in un’attività specifica. Serve a rendere esplicito un procedimento, ma non
	garantisce da solo la correttezza dell’esecuzione.
}

\voceglossario{societa-benefit}{Società Benefit}{
	Società che affianca allo scopo di lucro una o più finalità di beneficio comune indicate
	nello statuto. La disciplina italiana richiede una gestione che bilanci questi obiettivi e una
	relazione annuale sull’impatto prodotto; la qualifica non coincide con una certificazione privata
	di sostenibilità.
}

\voceglossario{software-engineer}{\emph{Software Engineer}}{
	Professionista che applica metodi di ingegneria del \emph{software} alla progettazione,
	realizzazione, verifica e manutenzione di sistemi \emph{software}.
}

\voceglossario{solution-technical-architect}{\emph{Solution \& Technical Architect}}{
	Figura che traduce le esigenze applicative in una soluzione tecnica, definendo
	l’architettura del sistema, le tecnologie e le integrazioni tra componenti.
}

\voceglossario{sprint}{\emph{Sprint}}{
	Ciclo di lavoro di durata fissa, non superiore a un mese nel riferimento Scrum, durante il
	quale il gruppo persegue un obiettivo concordato. La revisione dei risultati e la retrospettiva
	permettono di adattare il prodotto e il modo di lavorare al ciclo successivo.
}

\voceglossario{sse}{SSE}{
	\emph{Server-Sent Events}: standard del \emph{web} con cui un \emph{server} mantiene aperta la
	risposta a una richiesta HTTP e vi invia una sequenza di eventi testuali. È unidirezionale, dal
	\emph{server} al \emph{client}, e attraversa i comuni \emph{proxy}, purché non accumulino la
	risposta prima di inoltrarla.
}

\voceglossario{staging}{\emph{Staging}}{
	Ambiente separato dalla produzione in cui collaudare una versione dell’applicazione prima di
	renderla disponibile agli utenti. Riproduce le condizioni di esercizio per quanto necessario alla
	verifica, senza coincidere con il servizio effettivamente usato dai clienti.
}

\voceglossario{stand-up}{\emph{Stand-up}}{
	Breve riunione quotidiana, nata nei metodi agili, in cui ogni membro del gruppo riferisce su
	cosa ha lavorato, su cosa lavorerà e quali ostacoli incontra. Il nome viene dall’abitudine di
	tenerla in piedi, per mantenerla breve.
}

\voceglossario{startup}{\emph{Startup}}{
	Impresa giovane, di solito piccola, nata per sviluppare un prodotto o un servizio innovativo e
	per verificarne in fretta l’interesse sul mercato, con cicli di sviluppo brevi e ruoli poco
	formalizzati.
}

\voceglossario{storefront}{\emph{Storefront}}{
	Interfaccia di un negozio digitale rivolta al cliente, attraverso cui consultare il
	catalogo, selezionare prodotti e avviare l’acquisto. In Shopify può basarsi sui temi della
	piattaforma oppure su un’applicazione separata con architettura \emph{headless}.
}

\voceglossario{story-point}{\emph{Story point}}{
	Unità di stima relativa che un gruppo usa per confrontare le attività in termini di
	complessità, impegno e incertezza. Non corrisponde a un numero fisso di ore e acquista significato
	rispetto alle stime dello stesso gruppo; la scala di Fibonacci è una possibile convenzione, non un
	obbligo di Scrum.
}

\voceglossario{streaming}{\emph{Streaming}}{
	Trasmissione di una risposta in frammenti successivi, inviati man mano che vengono prodotti,
	così che chi la riceve possa mostrarla prima che sia completa. In un assistente basato su un
	modello linguistico permette di mostrare il testo mentre il modello lo genera.
}

\voceglossario{team-lead}{\emph{Team lead}}{
	Responsabilità di guida tecnica di un gruppo di lavoro. In Halue comprende l’analisi
	funzionale, le stime per il cliente e il coordinamento tecnico degli sviluppatori.
}

\voceglossario{theme-app-extension}{\emph{Theme app extension}}{
	Estensione con cui un’applicazione Shopify aggiunge elementi a un tema senza modificarne il
	codice. I file dell’estensione sono serviti dalla rete di Shopify e un solo rilascio raggiunge
	tutti i negozi che usano l’applicazione. Comprende blocchi che il commerciante posiziona in una
	sezione e blocchi \emph{app embed}, inseriti una volta in ogni pagina e attivati dall’editor del
	tema.
}

\voceglossario{ticket}{\emph{Ticket}}{
	Unità di lavoro registrata in un sistema di tracciamento, con un identificativo, una
	descrizione e uno stato. Può rappresentare una funzionalità, una modifica o un difetto, e
	raccogliere stime, discussioni e collegamenti alle modifiche del codice.
}

\voceglossario{token}{\emph{Token}}{
	Unità minima di testo elaborata da un modello linguistico: una parola breve, una parte di
	parola, un segno di punteggiatura. Il modello legge e genera testo un \emph{token} alla volta;
	la lunghezza del contesto e il costo di utilizzo si misurano in \emph{token}.
}

\voceglossario{token-accesso}{\emph{Token} di accesso}{
	Credenziale che un’applicazione presenta a ogni richiesta a un’API per dimostrare di essere
	autorizzata, con i permessi (\emph{scope}) che le sono stati concessi. Può essere permanente o
	avere una scadenza, nel qual caso si rinnova con un secondo \emph{token} (\emph{refresh token}).
	Da non confondere con il \emph{token} dei modelli linguistici.
}

\voceglossario{tool-calling}{\emph{Tool calling}}{
	Meccanismo con cui un modello linguistico chiede al programma che lo ospita di eseguire una
	funzione, detta strumento, indicandone il nome e gli argomenti. Il programma esegue la funzione e
	restituisce il risultato al modello, che lo usa per proseguire.
}

\voceglossario{tunnel}{\emph{Tunnel}}{
	Servizio che espone su un indirizzo pubblico un programma in
	esecuzione su una macchina non raggiungibile da Internet, inoltrandogli le richieste. Nello
	sviluppo di applicazioni Shopify serve perché Shopify possa chiamare l'applicazione eseguita sul
	computer dello sviluppatore.
}

\voceglossario{ucp}{UCP}{
	Protocollo aperto per rendere interoperabili i servizi commerciali e gli agenti che operano
	per conto dell’utente. Definisce modalità condivise per scoprire le capacità di un negozio e
	interagire con le fasi dell’acquisto; riguarda il dominio commerciale, mentre MCP offre un
	meccanismo generale per collegare applicazioni e strumenti.
}

\voceglossario{ux-ui-digital-consultant}{\emph{UX/UI Digital Consultant}}{
	Consulente che progetta l’esperienza utente e l’interfaccia di un prodotto digitale,
	considerando percorsi d’uso, struttura delle informazioni e aspetto visivo.
}

\voceglossario{widget}{\emph{Widget}}{
	Elemento dell’interfaccia autonomo, inserito in una pagina che non lo ha prodotto. Il
	\emph{widget} di ChatCart è il riquadro di chat sovrapposto alle pagine del negozio.
}

\end{description}
